\documentclass[12pt]{article}
\usepackage[utf8]{inputenc}
\usepackage[T1]{fontenc}
\usepackage{hyperref}
\title{Vaghua Language Resources for K-12 Curriculum Development}
\author{Compiled by Claude Sonnet 5 (claude-sonnet-5) for Andr\'as Kornai}
\date{}
\begin{document}
\maketitle

\section{Introduction}
Vaghua (ISO 639-3 tva; Glottocode vagh1249; also Vagua or Tavula) is an Austronesian language of Choiseul Island, Solomon Islands. It belongs to the West Choiseul group of Northwest Solomonic. Speaker counts differ by source: about 2,000 (Wikipedia, 1999) to about 4,000 (Joshua Project). Little documentation turned up in search.

\section{Primary Educational Resources}
None were found.

\section{Linguistic and Reference Materials}
\begin{itemize}
\item Wikipedia entry on Vaghua: \url{https://en.wikipedia.org/wiki/Vaghua_language}
\item Northwest Solomonic bibliography for Choiseul languages (Surrey Morphology Group): \url{https://www.smg.surrey.ac.uk/languages/northwest-solomonic/bibliography/choiseul/}
\item Lexicostatistical Study of the Languages of Choiseul Island, British Solomon Islands (ANU repository): \url{https://openresearch-repository.anu.edu.au/items/a48b5c1e-4825-4c41-8e12-53df64261c0c}
\end{itemize}

\section{Religious and Community Resources}
Joshua Project describes the Vagua Tavula people as mostly Christian and reports no published written form and Bible translation still needed: \url{https://joshuaproject.net/people_groups/15300/BP}

\section{Digital and Audio Resources}
None were found.

\section{Curriculum Development Context}
No orthography, dictionary, primer or school programme was found for Vaghua. Any curriculum would require new documentation work with the community, possibly drawing on the Choiseul bibliography above. This dossier is short because little exists in the search results.
\end{document}